\chapter*{Vorwort: Was passiert, wenn man die Exponenten umdreht?}
\addcontentsline{toc}{chapter}{Vorwort}
\markboth{VORWORT}{VORWORT}

Dieses Buch ist nicht am Schreibtisch entstanden, sondern im
Kinderzimmer, beim Zubettbringen meiner Tochter. Ich lag
neben ihr im Dunkeln und wollte nicht einschlafen, nur begleiten --
und nichts hält den Geist so zuverlässig wach wie ein Gedanke, der
ihn nicht loslässt. Bei mir sind das Zahlen; mein Kopf rechnet von
selbst. An jenem Abend übte ich Kopfrechnen, aus dem etwas
trotzigen Vorsatz, das Rechnen im Zeitalter der Computer nicht zu
verlernen: 728 geteilt durch 4. Die 700 sind 7 mal 100, also 7 mal
4 mal 25 -- durch 4 geteilt bleiben 7 mal 25, das ist 175; die 28
liefern trivial 7; zusammen 182.

Nichts daran ist bemerkenswert, außer dem, was mir beim Zerlegen
auffiel: wie selbstverständlich ich Hunderter, Zehner und Einer
auseinandersortiere. Von dort war es ein kurzer Weg zu dem, was ich
über Stellenwertsysteme wusste -- Binärzahlen, das Zwölfersystem,
das man an den Fingergliedern abzählen kann --, und zu einer Frage,
die ich so noch nirgends gestellt gesehen hatte: Was passiert, wenn
die Zählbasen \emph{innerhalb} einer Zahl wechseln? Dass es solche
Systeme längst gibt -- die Fakultätszahlen --, wusste ich in dieser
Nacht noch nicht; das ergab erst die erste Recherche. Aber
aufsteigende Basen kamen mir langweilig vor. Ich wollte es
andersherum: ganz rechts die meisten
Möglichkeiten.\footnote{Genau genommen wird kein Exponent
umgedreht, sondern die Richtung, in der die Basen wachsen -- die
Gewichtsordnung des Stellenwertsystems. Der Titel folgt dem
ursprünglichen Gedanken, nicht der Terminologie.} Wer diesen
Gedanken ernst nimmt, stößt sofort auf eine Zwangsläufigkeit: So
kann man nicht bis unendlich zählen -- Unendlichkeit wäre hier
Unsinn gewesen, also: endlich. Was zunächst wie ein Defekt
aussieht, ist die Geburt des zentralen Begriffs dieses Buchs: Jede
Zahl lebt in einem \emph{Horizont}, einem endlichen Raum mit einer
festen Anzahl von Plätzen. Meine Tochter schlief da längst. Meine Frau
erzählte mir später, ich sei gut eine Stunde oben gewesen; gefühlt
waren es fünfundzwanzig Minuten.

Dass führende Nullen in so einem System die Zahlen verändern -- die
Merkwürdigkeit, die alles Weitere antreibt --, wusste ich in jener
Nacht noch nicht. Sie kam ans Licht, wie man Ideen heute prüft: Ich
setzte mich an den Computer und erklärte den Gedanken GPT, mit der
Bitte, nachzusehen, ob es das alles schon gibt. Das erste Urteil
war trocken: Unsinn. Aber wir blieben dran, und nach einiger Zeit
merkten wir gemeinsam, dass es kein kompletter Unsinn war. In
diesen Gesprächen tauchten die Nullen auf, die vorne an einer Zahl
plötzlich etwas bedeuten, und viele der Gedanken, die dieses Buch
aufgreift. Ein Satz aus dieser Zeit ist mir geblieben: 4 mal 5 kann
eine Fläche sein -- und wenn die 4 ein Skalar ist, eine Länge; in
beiden Fällen 20. Mit der Horimetrik hat das nur am Rande zu tun,
aber es machte Eindruck -- und findet in
Kapitel~\ref{ch:strukturpolynome} ein spätes Echo. Der Name
übrigens entstand ganz unfeierlich: „Horizont-Arithmetik“ war mir
zu lang. Also Horimetrik.

Irgendwann war das Thema zu groß für Plaudereien, und GPT und ich
waren uns einig, dass es einen eigenen Agenten verdient. So kam
Fable ins Rennen -- das Claude-Modell von Anthropic, das auf der
Titelseite als mathematischer Mitarbeiter firmiert. Ich sage es
ehrlich: Mathematisch wurde ich schnell abgehängt. Mein
Informatikstudium hat mir Mathematik beigebracht, samt Nebenfach --
aber das liegt über ein Jahrzehnt zurück. Die Grundzüge verstehe
ich, und die waren spannend genug: nicht, weil es eine Anwendung
gäbe, sondern weil das alles neuartig genug wirkte, um ihm
nachzugehen. Viele mathematische Felder waren zu Beginn ohne
Anwendung. Fable kann Mathe; also ließ ich alles nachrechnen und
eine saubere Theorie aufstellen, und mit der Zeit fanden sich immer
mehr Übergänge zur etablierten Mathematik. GPT steuerte
zwischendurch weiter Ideen bei. Beiden gilt mein Dank -- lesen Sie
selbst, was die KI und ich gefunden haben.

Das Buch verfolgt zwei Ziele. Erstens will es die Mathematik dieses
Raums ehrlich entwickeln -- mit Definitionen, Beweisen und auch mit
den Irrwegen, die wir unterwegs gegangen sind. Zweitens will es
lesbar bleiben: Jede technische Passage kehrt an die Oberfläche
zurück, bevor die nächste beginnt. Wer nur die formelfreien
Abschnitte liest, soll trotzdem verstehen, worum es geht.

Ob die Horimetrik eine eigenständige Theorie wird oder eine
besonders lehrreiche Spielerei bleibt, wussten wir beim Schreiben
selbst noch nicht. Genau davon handelt dieses Buch.
